\documentclass[aspectratio=169]{beamer}
\usetheme{Madrid}
\usecolortheme{default}
\usepackage{amsmath}
\usepackage{graphicx}
\usepackage{tikz-cd}
\usepackage{multicol}

\setlength{\parindent}{0pt}
\setlength{\parskip}{1\baselineskip}

\title[Lecture 10]{An Introduction to E-Commerce\\\small with a Focus on Machine Learning, Deep Learning, and Artificial Intelligence}
\subtitle{Lecture 10: Fraud, Risk, and Trust \& Safety}
\author{Dr. Haitham A. El-Ghareeb}
\institute{Faculty of Computers and Information Sciences \\ Mansoura University \\ Egypt}
\date{\today}

\begin{document}

\begin{frame}
  \titlepage
\end{frame}

\begin{frame}{Session plan (90 minutes)}
  \begin{enumerate}
    \item Why fraud and trust \& safety matter
    \item Fraud and abuse categories
    \item Fraud as a machine learning problem
    \item Signals and features for fraud detection
    \item Rule-based systems
    \item Hands-on lab: risk scoring and thresholding
  \end{enumerate}
\end{frame}

\begin{frame}{Learning objectives}
  By the end of this lecture, you should be able to:
  \begin{itemize}
    \item Identify major categories of fraud and abuse in e-commerce, including payment fraud, account takeover, promotion abuse, and refund abuse.
    \item Explain why fraud detection is a cost-sensitive, adversarial, and highly imbalanced decision problem.
    \item Compare supervised learning, anomaly detection, graph-based methods, and rule+ML hybrid approaches.
    \item Design a trust-and-safety workflow that combines scoring, thresholds, review operations, and business guardrails.
  \end{itemize}
\end{frame}

\section{Why fraud and trust \& safety matter}

\begin{frame}{Why fraud and trust \& safety matter}
  \begin{itemize}
    \item payment costs,
    \item customer trust,
    \item seller trust in marketplaces,
    \item operational review burden,
    \item promotional efficiency,
    \item regulatory and compliance exposure.
  \end{itemize}
\end{frame}

\section{Fraud and abuse categories}

\begin{frame}{Fraud and abuse categories}
  \begin{itemize}
    \item \textbf{Payment fraud:} Payment fraud often involves stolen cards, synthetic identities, mule accounts, or coordinated abuse of payment flows.
    \item \textbf{Account takeover:} Account takeover occurs when an attacker gains control of an existing customer account.
    \item \textbf{Promotion and coupon abuse:} creating multiple accounts to redeem new-user offers,; coordinated coupon sharing,
    \item \textbf{Refund and return abuse:} Refund abuse includes false non-delivery claims, empty-box returns, repeated wardrobing behavior, and manipulation of return policies.
    \item \textbf{Marketplace and seller abuse:} counterfeit listings,; fake reviews,
  \end{itemize}
\end{frame}

\section{Fraud as a machine learning problem}

\begin{frame}{Fraud as a machine learning problem}
  \begin{itemize}
    \item \textbf{Class imbalance:} Fraudulent events are typically a small fraction of all transactions.
    \item \textbf{Adversarial behavior:} Once a platform changes a rule or scoring model, attackers may alter devices, payment instruments, shipping behavior, or timing patterns.
    \item \textbf{Label delay and label noise:} Ground truth often arrives late.
    \item \textbf{Asymmetric costs:} False negatives lose money and trust.
  \end{itemize}
\end{frame}

\section{Signals and features for fraud detection}

\begin{frame}{Signals and features for fraud detection}
  \begin{itemize}
    \item \textbf{Transaction signals:} order value,; payment method,
    \item \textbf{Account signals:} account age,; prior successful orders,
    \item \textbf{Device and network signals:} device fingerprint,; IP reputation,
    \item \textbf{Behavioral and sequence signals:} login and checkout timing,; session depth,
  \end{itemize}
\end{frame}

\section{Rule-based systems}

\begin{frame}{Rule-based systems}
  \begin{itemize}
    \item \textbf{Examples of useful rules:} block transactions from known bad cards or IPs,; trigger review if order value exceeds a threshold and account age is very low,
    \item \textbf{Strengths and weaknesses:} deterministic policy enforcement,; emergency response,
  \end{itemize}
\end{frame}

\section{Supervised learning}

\begin{frame}{Supervised learning}
  \begin{itemize}
    \item \textbf{Typical model families:} logistic regression,; tree-based ensembles,
    \item \textbf{Why supervised learning helps:} aggregating many weak predictors,; ranking risk more smoothly than rigid rules,
    \item \textbf{Operational caution:} the quality of labels,; the recency of training data,
  \end{itemize}
\end{frame}

\section{Anomaly detection}

\begin{frame}{Anomaly detection}
  \begin{itemize}
    \item \textbf{What anomaly methods look for:} unusual transaction timing,; atypical device usage,
    \item \textbf{Strengths and weaknesses:} Anomaly methods are useful for surfacing suspicious patterns and complementing supervised systems.
  \end{itemize}
\end{frame}

\section{Graph-based fraud analysis}

\begin{frame}{Graph-based fraud analysis}
  \begin{itemize}
    \item \textbf{Why graph structure matters:} many accounts linked to one device,; many cards shipping to related addresses,
    \item \textbf{Graph features and graph ML:} connected-component and degree features,; shared-entity counts,
  \end{itemize}
\end{frame}

\section{Hybrid rules + ML systems}

\begin{frame}{Hybrid rules + ML systems}
  \begin{itemize}
    \item \textbf{Why hybrids work:} enforce mandatory blocks,; assign risk scores,
    \item \textbf{Example decision flow:} apply deterministic policy rules,; compute risk score from supervised and anomaly signals,
  \end{itemize}
\end{frame}

\section{Thresholds and decision economics}

\begin{frame}{Thresholds and decision economics}
  \begin{itemize}
    \item \textbf{Approve, review, reject:} \textbf{Approve:} low-risk transactions proceed automatically.; \textbf{Review:} medium-risk cases go to human review or secondary checks.
    \item \textbf{Cost-sensitive tuning:} fraud loss,; customer lifetime value,
  \end{itemize}
\end{frame}

\section{Human-in-the-loop review}

\begin{frame}{Human-in-the-loop review}
  \begin{itemize}
    \item \textbf{What reviewers need:} transaction history,; device and network context,
    \item \textbf{Feedback loop:} rule refinement,; model retraining,
  \end{itemize}
\end{frame}

\section{Fraud evaluation metrics}

\begin{frame}{Fraud evaluation metrics}
  \begin{itemize}
    \item \textbf{Useful metrics:} precision,; recall,
    \item \textbf{Business interpretation:} A model with strong recall but excessive false positives may block good customers.
  \end{itemize}
\end{frame}

\section{Drift, adaptation, and monitoring}

\begin{frame}{Drift, adaptation, and monitoring}
  \begin{itemize}
    \item \textbf{What to monitor:} score distribution shifts,; feature drift,
    \item \textbf{Why monitoring is essential:} fraud patterns shift,; product flows change,
  \end{itemize}
\end{frame}

\section{Trust \& safety beyond payment fraud}

\begin{frame}{Trust \& safety beyond payment fraud}
  \begin{itemize}
    \item \textbf{Examples beyond transactions:} review authenticity,; marketplace seller compliance,
    \item \textbf{Why this matters architecturally:} commerce events,; ERP refund records,
  \end{itemize}
\end{frame}

\section{Hands-on lab: risk scoring and thresholding}

\begin{frame}{Hands-on lab: risk scoring and thresholding}
  \begin{itemize}
    \item \textbf{Lab scenario:} Students are given a synthetic transaction dataset containing payment, account, device, and label information, plus estimated fraud-loss and false-positive costs.
    \item \textbf{Lab tasks:} Build a baseline risk model using structured transaction and account features.; Compute precision, recall, and a cost-sensitive evaluation summary.
    \item \textbf{Expected learning outcome:} By the end of the lab, students should be able to explain fraud detection as a staged decision system rather than a single prediction score.
  \end{itemize}
\end{frame}

\section{Managerial implications}

\begin{frame}{Managerial implications}
  \begin{itemize}
    \item clear policy,
    \item strong data quality,
    \item hybrid detection methods,
    \item review operations,
    \item and continuous monitoring.
  \end{itemize}
\end{frame}

\end{document}
